\newif\ifglobal

%%%%%%%%%%%%%%%%%%%%%%%
%%% Compile to view %%%
%%%%%%%%%%%%%%%%%%%%%%%

%%%%%%%%%%%%%%%%%%%%%%%%%%%%%%%%%%%%%%%%%%%%%%%%%%%
%%% Readme:                                     %%%
%%% https://www.overleaf.com/read/bwdjvfjksvjy/ %%%
%%%%%%%%%%%%%%%%%%%%%%%%%%%%%%%%%%%%%%%%%%%%%%%%%%%

% >>> M?: marks a module setting number or important notice

% >>> M4: Set {./relative-path} to external files for this module <<<
\def \modulefiles {./01-RISCVCPU}
% To add an external file, use:
% (Replace '---' with actual filename)
% '\input{\modulefiles/tables/---.tex}' for tables
% '\includegraphics{\modulefiles/figures/---}' for figures
% '\subfile{\modulefiles/---}' for register files

% >>> M1: This module is in global project? <<<
% 'YES' - uncomment; 'NO' - comment out
\globaltrue

\ifglobal
    \documentclass[main\_\_EN.tex]{subfiles}

\else
    \documentclass[openany, 10pt]{book}

    % >>> M2: In ./00-shared/preamble-trm-module, also find
    %         settings for visibility of todo notes, watermark <<<
    \usepackage{preamble-trm-module}

    % >>> M3: Temporary labels in English,
    %         you can add your own labels there <<<
    \myexternaldocument{./00-shared/temp-labels-en}

\fi %\ifglobal


\setcounter{secnumdepth}{4}


% Define formatting of parts and chapters in text
\titleformat{\chapter}
  [display]
  {\LARGE\bfseries}
  {\color{AllportsColor}Chapter \thechapter}
  {1em}
  {\color{AllportsColor}}
\titlespacing*{\chapter}{0pt}{0pt}{20pt}

% Define headers
\usepackage{fancyhdr}
\usepackage{lipsum}
\renewcommand{\chaptermark}[1]{\markboth{Chapter\, \thechapter\, \ #1}{}}
\fancypagestyle{plain}{%
  \fancyhead{}%
  %\fancyhead[R]{\Navigation}
  \fancyhead[L]{%
    \nouppercase{%
      \ %
      \normalsize\leftmark
    }%
  }%
}
\pagestyle{plain}


% >>> M5: If this module needs any updates in
% './00-shared/preamble-trm-module.sty'
% (include additional package, etc.), contact your module PM
% or create the file './00-shared/changelog.tex' make the
% necessary updates yourself and document them in this file <<<

% Make text left-aligned and allow latex to break pages naturally, comment out for full justification
\raggedright
\raggedbottom


\begin{document}


% Sets language into which pre-defined bits of text are translated
\selectlanguage{english}

% Do not delete this block
\ifglobal
% Add the following front matter if
% this module is outside of global project
\else
    \InputIfFileExists{readme}{}{}
    \listoftodos
    {\let\clearpage\relax \listofdonetodos}
    \tableofcontents
    %\listoftables
    %\listoffigures
\fi


%%%%%%%%%%%%%%%%%%%%%%%%%
%%% TRM Module Begins %%%
%%%%%%%%%%%%%%%%%%%%%%%%%

\hypertarget{riscvcpu}{}
\chapter{ESP-RISC-V CPU}
\label{mod:riscvcpu}





\section{Overview}
\chipname{} is powered by Espressif's second-generation 32-bit CPU core complex based upon RISC-V Instruction Set Architecture (ISA). The CPU Core complex consists of a RISC-V CPU core with a dedicated core-local interrupt controller (CLIC), a debug block, and a core-local interrupt (CLINT) timer. Figure \ref{fig:riscvcpu-core-complex-block-diagram} shows the block diagram of the CPU core complex.

\begin{figure}[H]
\centering
\includegraphics[scale=1,width=0.8\textwidth]{\modulefiles/figures/C61_block_diag.png}
\caption{CPU Core Complex Block Diagram}
\end{figure}
\label{fig:riscvcpu-core-complex-block-diagram}

RISC-V core in CPU core complex uses a 5-stage, in-order, scalar pipeline architecture optimized for area, power, and performance. RISC-V core also has built-in dedicated core-local interrupt controller (CLIC) and system bus (SYS BUS) interfaces for memory and peripheral access. The DEBUG block provides a connection between the JTAG interface and the RISC-V core. The core also provides an instruction trace interface for offline debugging.

\section{Features}
\label{sec:riscvcpu-features}
RISC-V HP CPU core in the CPU core complex implements the following standard RISC-V extensions:
\begin{itemize}
    \item Mandatory base integer (I)
    \item Multiplication/Division (M)
    \item Atomic (A)
    \item Compressed (C)
    \item Code size reduction (Zc)
    \item Bit manipulation (Zb)
\end{itemize}

Apart from the above RISC-V instruction extensions support, CPU core supports the features below:
\begin{itemize}
    \item Five-stage pipeline that supports an operating clock frequency up to 160 MHz
    \item Compatible with RISC-V Instruction Set Manual, Volume I: RISC-V User-Level ISA, Version 2.2
    \item Compatible with RISC-V Instruction Set Manual, Volume II: RISC-V Privileged Architecture, Version 1.10
    \item Compatible with RISC-V Core-Local Interrupt Controller (CLIC) Version 0.9
    \item Compatible with RISC-V External Debug Support Version 0.13.2
    \item Zero wait cycle access to on-chip SRAM and cache for program and data access over IRAM/DRAM interface
    \item 32-bit AHB system bus for peripheral access
    \item Core-local interrupts (CLINT) support
    \item User (U) privilege mode execution
    \item User (U) mode interrupt delegation via CLIC
    \item Bus error response
    \item Standard physical memory protection (PMP) and custom attributes (PMA) for up to 16 configurable regions for security
    \item Branch target buffer (BTB) with BHT and RAS for performance improvement
    \item Debug module (DM) compliant with the specification RISC-V External Debug Support Version 0.13.2 with external debugger support over an industry-standard JTAG/USB port
    \item Debugger with a direct system bus access (SBA) to memory and peripherals via the core data bus
    \item Support for instruction trace for offline debug
    \item Hardware trigger compliant with the specification RISC-V External Debug Support Version 0.13.2 with up to 3 breakpoints/watchpoints
    \item Support for up to 8 dedicated IOs
    \item Configurable events for core performance metrics
\end{itemize}

\section{Terminology}

To better illustrate the functionality of ESP-RISC-V CPU, the following terms are used in this chapter.

\begin{longtable}[c]{ R{2cm} p{12cm} }
\textbf{ABI} & Application binary interface\\
\textbf{branch} & An instruction that conditionally changes the execution flow \\
\textbf{CLIC} & Core-local interrupt controller\\
\textbf{CLINT} & Core-local interrupt\\
\textbf{CSR} & Control and status register\\
\textbf{delta} & A change in the program counter that is other than the difference between two instructions placed consecutively in memory \\
\textbf{DTM} & Debug transport module \\
\textbf{FPR} & Floating-point register\\
\textbf{GPIO} & General-purpose input/output\\
\textbf{GPR} & General-purpose register\\
\textbf{HWLP} & Hardware loop\\
\textbf{hart} & RISC-V hardware thread \\
\textbf{LSB} & Least significant bit\\
\textbf{LR} & Load reserved\\
\textbf{MSB} & Most significant bit\\
\textbf{PMP} & Physical memory protection\\
\textbf{retire} & The final stage of executing an instruction, when the machine state is updated \\
\textbf{RO} & Read-only field\\
\textbf{R/W} & Readable and writable field\\
\textbf{SC} & Store conditional\\
\textbf{SoC} & System on chip\\
\textbf{WO} & Write-only field\\
\textbf{WARL} & Write any read legal\\
\textbf{W1} & Write One\\
\end{longtable}

\section{Address Map}
The table below shows the address map of various regions accessible by CPU for instruction, data, and system bus peripherals.

\begin{longtable}{| l| l |c | c | c  |}
\caption{CPU Address Map} \label{tab:riscvcpu-mem-blocks}
\\ \hline
\rowcolor{lightgray}
\textbf{Region} & \textbf{Description}    & \textbf{Start Address}  & \textbf{End Address}  & \textbf{Access Type}   \\ \hline
CPU \label{cpu} &  CLIC/CLINT registers &      0x2000\_0000  &  0x2FFF\_FFFF  &  R/W  \\ \hline
IRAM/DRAM\label{ram} & Instruction/Data region &   0x4000\_0000  &  0x4FFF\_FFFF  &  R/W  \\ \hline
CPU peripherals \label{cpu_peri} & AHB (Strong order access)  &      0x6000\_0000  &  0x6FFF\_FFFF  &  R/W  \\ \hline
AHB\label{sysahb} & AHB (Weak Order Access) &      *Default  &  *Default  &  R/W  \\ \hline
\end{longtable}

\textbf{*Default:} Address not matching any of the specified ranges above are accessed through AHB bus by a CPU core.

%%%%%%%%%%%%%%%%%%%%%%%%%%%
%%% C61 Register Summary %%%
%%%%%%%%%%%%%%%%%%%%%%%%%%%

\section{Configuration and Status Registers (CSRs)}
\hypertarget{riscvcpu-cpu-reg-summ}{}
\subsection{Register Summary}
\label{sec:riscvcpu-cpu-reg-summ}

Below is a list of CSRs implemented for HP CPU core. All the implemented CSRs follow the standard mapping of bit fields as described in the RISC-V Instruction Set Manual, Volume II: Privileged Architecture, Version 1.10. It must be noted that even among the standard CSRs, not all bit fields have been implemented, limited by the subset of features implemented in the CPU. Refer to the next section for detailed description of the subset of fields implemented under each of these CSRs.

\subfile{\modulefiles/01-RISCVCPU-reg-summ__EN}
Note that if write/set/clear operation is attempted on any of the CSRs which are read-only (RO), as indicated in the above table, the CPU will generate an illegal instruction exception.

\subsection{Register Description}

\subfile{\modulefiles/01-RISCVCPU-reg__EN}











%%%%%%%%%%%%%%%%%%%%%%%%
%%% RISC-V Standard Extensions %%%
%%%%%%%%%%%%%%%%%%%%%%%%
\newpage
\section{RISC-V Standard ISA Extensions Support}


HP CPU core implements mandatory base integer ISA and is compliant with version 2.0 of RV32I base integer instruction set specified in RISC-V Instruction Set Manual, Volume I: RISC-V User-Level ISA, Version 2.2

%%%%%%%%%%%%%%%%%%%%%%%%
%%% M Extension %%%
%%%%%%%%%%%%%%%%%%%%%%%%
\subsection{Multiplication/Division (M) Extension}
\label{sec:riscvcpu-mext}
\subsubsection{Overview}
\label{sec:riscvcpu-mext-overview}
In addition to base RV32I instruction set, HP core also implements standard integer multiplication and division extension and is fully compliant with version 2.0 of standard extension "M" specified RISC-V Instruction Set Manual, Volume I: RISC-V User-Level ISA, Version 2.2.

\subsubsection{Functional Description}
\label{sec:riscvcpu-mext-func-descr}

Instructions specified in "M" extension perform multiply and divide operations on the values held in two integer registers. Below are CPU cycle details taken by multiply and divide instructions in the absence of any data hazard in the pipeline.

\begin{itemize}
    \item MUL operation takes 2 cycles
    \item DIV operation takes 1-19 cycles depending on operands
\end{itemize}


%%%%%%%%%%%%%%%%%%%%%%%%
%%% C Extension %%%
%%%%%%%%%%%%%%%%%%%%%%%%
\subsection{Compressed (C) Extension}
\label{sec:riscvcpu-cext}

\subsubsection{Overview}
\label{sec:riscvcpu-cext-overview}

HP core also implements RISC-V standard compressed instruction set extension named "C" and is fully compliant with version 2.0 of standard extension "C" specified in RISC-V Instruction Set Manual, Volume I: RISC-V User-Level ISA, Version 2.2. RVC is generic term used to specify C extension support. RVC reduces static and dynamic code size by adding short 16-bit instruction encoding for common operations. Typically, 50-60\% of the RISC-V instructions in a program can be replaced with RVC instructions, resulting in a 25-35\% code-size reduction.

\subsubsection{Functional Description}
\label{sec:riscvcpu-cext-func-descr}
RVC uses a simple compression scheme that offers a shorter 16-bit version of 32-bit RISC-V instructions for any of the below scenarios:
\begin{enumerate}
    \item Immediate and address offset is small
    \item One of the registers is a zero register (x0), the ABI link register (x1), or the ABI stack pointer (x2)
    \item Destination register and the first source register are identical, or
    \item The registers used are the 8 most popular ones i.e. x8-x15.
\end{enumerate}

\newpage

%%%%%%%%%%%%%%%%%%%%%%%%
%%% Zc Extension %%%
%%%%%%%%%%%%%%%%%%%%%%%%

\newpage
\subsection{Code Size Reduction (Zc) Extension}
\label{sec:riscvcpu-zc}

\subsubsection{Overview}
\label{sec:riscvcpu-zc-overview}

Zc extension is useful for code size reduction. HP core implements this extension compatible with \href{https://github.com/riscvarchive/riscv-code-size-reduction/blob/v1.0.4-3/Zc-specification/Zc.adoc}{Zc-specification v1.0.4-3}. The Zc-specification has multiple subset extensions, and the \chipname HP CPU has implemented Zcb, Zcmp, and Zcmt extensions.

\subsubsection{Functional Description}
\label{sec:riscvcpu-zc-func-descr}








\subparagraph{Zcmt Extension}

The Zcmt extension is also referred to as table jump. It reads the jump target address from the jump table and then jumps to it.

Table jump uses a 256-entry 32-bit wide table in instruction memory to contain function address. The table must be a minimum of 64-byte aligned. It is used as a form of dictionary compression to reduce the code size of jal, auipc+jalr, jr, auipc+jr instructions.

Table jump allows the linker to replace the following instruction sequences with a cm.jalt encoding, and an entry in the table:
\begin{itemize}
    \item 32-bit j calls
    \item 32-bit jal ra calls
\end{itemize}

\subsubsection{Zc Instructions}
\label{sec:riscvcpu-zc-insn-descr}

\begin{longtable}{ | p{2.0cm} | p{4.0cm} |  p{10.0cm} | }
\hline\rowcolor{lightgray}
Instruction & Mnemonic & Description  \\ \hline
\endfirsthead
\hline
\rowcolor{lightgray}
Instruction & Mnemonic & Description  \\ \hline
\hline
\endhead
\endfoot

\multicolumn{3}{|l|}{\textbf{Zcb extension}} \\\hline
c.lbu & c.lbu rd', uimm(rs1') & Loads a byte from the memory address formed by adding rs1' to the zero-extended immediate uimm. The resulting byte is zero-extended to 32 bits and is written to rd' \\\hline
c.lhu & c.lhu rd', uimm(rs1') &  Loads a halfword from the memory address formed by adding rs1' to the zero-extended immediate uimm. The resulting byte is zero-extended to 32 bits and is written to rd'\\\hline
c.lh & c.lh rd', uimm(rs1') & Loads a halfword from the memory address formed by adding rs1' to the zero-extended immediate uimm. The resulting byte is sign-extended to 32 bits and is written to rd' \\\hline
c.sb & c.sb rs2', uimm(rs1') & Stores the least significant byte of rs2' to the memory address formed by adding rs1' to the zero-extended immediate uimm \\\hline
c.sh & c.sh rs2', uimm(rs1') &  Stores the least significant halfword of rs2' to the memory address formed by adding rs1' to the zero-extended immediate uimm \\\hline
c.zext.b & c.zext.b rsd' & Zero-extends the least significant byte of the operand to 32 bits by inserting zeros into all of the bits more significant than 7 \\\hline
c.sext.b & c.sext.b rsd' & Sign-extends the least significant byte of the operand to 32 bits by copying the most significant bit in the byte to all the more significant bits   \\\hline
c.zext.h & c.zext.h rsd' &  Zero-extends the least significant halfword of the operand to 32 bits by inserting zeros into all of the bits more significant than 15 \\\hline
c.sext.h  & c.sext.h rsd'  &  Sign-extends the least significant halfword of the operand to 32 bits by copying the most significant bit in the halfword to all the more significant bits \\\hline
c.not & c.not rsd' & Takes the one's complement of rd'/rs1' and writes the result to the same register \\\hline
c.mul & c.mul rsd' & Multiplies 32 bits of the source operands from rsd' and rs2' and writes the lowest 32 bits to the result to rsd'\\\hline

\multicolumn{3}{|l|}{\textbf{Zcmp extension}} \\\hline

cm.push & cm.push \{reg\_list\}, -stack\_adj & Stores the registers in reg\_list to the memory below the stack pointer, and then creates the stack frame by decrementing the stack pointer by stack\_adj, including any additional stack space requested by the value of spimm\\\hline
cm.pop & cm.pop \{reg\_list\}, stack\_adj & Loads the registers in reg\_list from stack memory, and then adjusts the stack pointer by stack\_adj \\\hline
cm.popret & cm.popret \{reg\_list\}, stack\_adj & Loads the registers in reg\_list from stack memory, adjusts the stack pointer by stack\_adj, then return to ra\\\hline
cm.popretz & cm.popretz \{reg\_list\}, stack\_adj & Loads the registers in reg\_list from stack memory, adjusts the stack pointer by stack\_adj, moves zero to a0 and then return to ra\\\hline
cm.mva01s & cm.mva01s rs1', rs2' &  Moves r1s' into a0 and r2s' into a1\\\hline
cm.mva01 & cm.mva01 r1s', r2s' & Moves a0 into r1s' and a1 into r2s' \\\hline

\multicolumn{3}{|l|}{\textbf{Zcmt extension}} \\\hline

cm.jt & cm.jt index & Reads an entry from the jump vector table in memory and jumps to the address that was read \\\hline
cm.jalt & cm.jalt index & Reads an entry from the jump vector table in memory and jumps to the address that was read, linking to ra  \\\hline

\end{longtable}

For more details about the instructions please refer to \href{https://github.com/riscvarchive/riscv-code-size-reduction/blob/v1.0.4-3/Zc-specification/Zc.adoc}{Zc-specification v1.0.4-3}.

\subsubsection{Limitations}
\label{sec:riscvcpu-zc-limit-descr}

\hyperref[regdesc:JVT]{jvt} adds architecture state to the context. Therefore, it must be saved and restored on context switches.

%%%%%%%%%%%%%%%%%%%%%%%%
%%% Zb Extension %%%
%%%%%%%%%%%%%%%%%%%%%%%%

\newpage
\subsection{Bit Manipulation (Zb) Extension}
\label{sec:riscvcpu-zb}

\subsubsection{Overview}
\label{sec:riscvcpu-zb-overview}

The bit-manipulation (bitmanip) extension, also known as the Zb extension set, consists of several component extensions to the base RISC-V architecture. These extensions aim to improve code density, performance, and energy efficiency.  The HP core fully supports the Zb extension. For Zb extension specifications, please refer to \href{https://github.com/riscv/riscv-b}{Zb-specification v1.0.0}.

\subsubsection{Functional Description}
\label{sec:riscvcpu-zb-func-descr}

Although the bitmanip instructions are designed for general-purpose use, certain instructions are more applicable to specific domains. Hence, the bitmanip extension set is divided into several smaller extensions instead of a single large one. Each extension has its own Zb*-extension name, and comprises instructions that share related functionality, use cases, and often the same implementation logic. Some instructions are available in only one extension, while others are available in several. The instructions have mnemonics and encodings that are independent of the extensions in which they appear. Thus, when implementing extensions with overlapping instructions, there is no redundancy in logic or encoding.

\newpage

%%%%%%%%%%%%%%%%%%%%%%%%
%%% Atomic Extension %%%
%%%%%%%%%%%%%%%%%%%%%%%%

\newpage
\subsection{Atomic (A) Extension}
\label{sec:riscvcpu-atomic}

\subsubsection{Overview}
\label{sec:riscvcpu-atomic-overview}

Support for atomic (A) extension is available in compliance with RISC-V Instruction Set Manual, Volume I: RISC-V User-Level ISA, Version 2.2, with an emphasis to guarantee forward progress, i.e. any situation that may cause data memory lock for an indefinite amount of time is prevented by their very functionality.

\subsubsection{Functional Description}
\label{sec:riscvcpu-atomic-func-descr}

Atomic (A) Extension contains instructions that atomically read-modify-write memory to support synchronization between multiple RISC-V harts running in the same memory space. The two forms of atomic instruction provided are load-reserved/store-conditional instructions and atomic fetch-and-op memory instructions. The atomic instructions currently ignore the aq (acquire) and rl (release) bits as they are irrelevant to the current architecture, in which memory ordering is always guaranteed.

\subsubsection{Load-Reserve (LR.W) Instruction}
\label{sec:riscvcpu-atomic-lr}

The LR.W instruction simply locks a 32-bit aligned memory address to which the load access is being performed. Once a 4-byte memory region is locked, it will remain locked, i.e. other harts won't be able to access this same memory location, until any of the following scenarios is encountered during execution:

\begin{itemize}
    \item any load operation
    \item any store operation
    \item any interrupts/exceptions
    \item backward jump/taken backward branch
    \item JALR
    \item ECALL/EBREAK/MRET/URET
    \item FENCE/FENCE.I
    \item debug mode
    \item critical section exceeding 64 bytes
    \item data address in SC.W instruction not matching that in LR.W instruction
\end{itemize}

If any of the above happens, except SC.W, the memory lock will be released immediately. If an SC instruction is encountered instead, the lock will be released eventually (not immediately) in the manner described in Section \ref{sec:riscvcpu-atomic-sc}.

If a misaligned address is encountered, it will cause an exception with \hyperref[regdesc:MCAUSE]{mcause} = 6.

\subsubsection{Store-Conditional (SC.W) Instruction}
\label{sec:riscvcpu-atomic-sc}

The SC.W instruction first checks if the memory lock is still valid, and the address is the same as specified during the last LR.W instruction. If so, only then will it perform the store to memory, and later release the lock as soon as it gets an acknowledgement of operation completion from the memory.

On the other hand, if the lock is found to have been invalidated (due to any of the situations as described in Section \ref{sec:riscvcpu-atomic-lr}), it will set a fail code (currently always 1) in the destination register rd.

If a misaligned address is encountered, it will cause an exception with \hyperref[regdesc:MCAUSE]{mcause} = 6.

\subsubsection{AMO Instructions}
\label{sec:riscvcpu-atomic-amo}

An AMO instruction executes in three steps:

\begin{enumerate}
    \item Read data from memory address given by rs1, and save it to destination register rd.
    \item Combine the data in rd and rs2 according to the operation type and keep the result for Step 3 below.
    \item Write the result obtained in Step 2 above to memory address given by rs1.
\end{enumerate}

There are nine different AMO operations: SWAP, ADD, AND, OR, XOR, MAX, MIN, MAXU, and MINU.

During this whole process, the memory address is kept locked from being accessed by other harts. If a misaligned address is encountered, it will cause an exception with \hyperref[regdesc:MCAUSE]{mcause} = 6.

For AMO operations both load and store access faults (PMP/PMA) are checked in the 1st step itself. For such cases \hyperref[regdesc:MCAUSE]{mcause} = 7.

%%%%%%%%%%%%%%%%%%%%%%%%
%%% Memory-Mapped Registers %%%
%%%%%%%%%%%%%%%%%%%%%%%%
\newpage
\section{Memory-Mapped Registers}
\label{sec:riscvcpu-tcip}

\subsection{Overview}
\label{sec:riscvcpu-tcip-overview}


The tightly-coupled memory-mapped registers below are present within the core complex:

\begin{itemize}
    \item CLIC registers
    \item CLINT registers
\end{itemize}

\subsection{Features}
\label{sec:riscvcpu-tcip-features}
HP CPU core supports:
\begin{itemize}
    \item Access to CLIC registers
    \item Bus-error exception generation on invalid address access
    \item Access to CLINT registers
\end{itemize}

\subsection{Functional Description}
\label{sec:riscvcpu-tcip-func-descr}

The table below shows the address decoding used by HP core to access CLIC and CLINT address range.

\begin{longtable}{ | l | p{5 cm} | }
\caption{Address Decoding for Memory-Mapped Register Access} \label{tab:riscvcpu-tcip-table}
\\ \hline
\rowcolor{lightgray}
Bits & Description \\ \hline
31:28 & Fixed as 0x2, i.e. base address \\ \hline
27:20 & Identifies the sub module: \newline 0x0: CLINT \newline 0x8: CLIC \newline 0xB: CLIC user \newline Others: Reserved \\ \hline
19:16 & Controls access type \newline 0x0: Access own registers \newline 0x1: Reserved\\ \hline
15:0 & Register offset \\ \hline
\end{longtable}


\begin{longtable}{ | l | c | c |}
\caption{Internal Memory Address Map} \label{tab:riscvcpu-tcip-slave-table}
\\ \hline
\rowcolor{lightgray}
Block & Start Address & End Address \\ \hline
CLINT & 0x20000000  &  0x2000FFFF\\ \hline
CLIC & 0x20800000  &  0x2080FFFF\\ \hline
\end{longtable}

%%%%%%%%%%%%%%%%%%%%%%%%%%%%
%%% Interrupt Controller %%%
%%%%%%%%%%%%%%%%%%%%%%%%%%%%
\newpage
\section{Interrupt Controller}
\label{sec:riscvcpu-int-ctrl}

%%%%%%%%%%%%%%%%%%
%%% INTERRUPTS %%%
%%%%%%%%%%%%%%%%%%

\subsection{Overview}
\label{sec:riscvcpu-intr-overview}


HP core uses the CLIC + CLINT scheme for implementing and supporting interrupts.

CLINT stands for core-local interrupts, which are sources local to the HP core subsystem for generating timer and software interrupts for HP core.

CLIC stands for core-local interrupt controller, which provides low-latency, vectored, pre-emptive interrupts for HP core.
The HP core CLIC unit supports 32 external interrupts and 2 CLINT interrupts.

%%%%%%%%%%%%%
%%% CLIC %%%
%%%%%%%%%%%%%

\subsection{CLIC}
\label{sec:riscvcpu-clic}

\subsubsection{Overview}
\label{sec:riscvcpu-clic-overview}


The HP core CLIC implementation follows the official specification: \href{https://github.com/riscv/riscv-fast-interrupt/blob/v0.9/src/clic.adoc}{Core-Local Interrupt Controller (CLIC) RISC-V Privileged Architecture Extensions (10/10/2023)}.


Only the CLIC mode of operation is supported by the HP core, i.e. \hyperref[fielddesc:MTVECMODE]{mtvec.MODE} is hardwired to 0x3. Hence, the basic RISC-V interrupt handling scheme and the associated CSRs (such as mie, mip, mideleg, uie, and uip) are unavailable.

HP core supports machine mode interrupt as well as user mode interrupt delegation.

\subsubsection{CLIC CSRs}
\label{sec:riscvcpu-clic-csrs}

The following machine mode CSRs are relevant for interrupt handling in the CLIC scheme:

\begin{multicols}{4}
\begin{itemize}
\item \hyperref[regdesc:MSTATUS]{mstatus}
\item \hyperref[regdesc:MTVEC]{mtvec}
\item \hyperref[regdesc:MTVT]{mtvt}
\item \hyperref[regdesc:MSCRATCH]{mscratch}
\item \hyperref[regdesc:MEPC]{mepc}
\item \hyperref[regdesc:MCAUSE]{mcause}
\item \hyperref[regdesc:MNXTI_REG]{mnxti}
\item \hyperref[regdesc:MINTTHRESH]{mintthresh}
\item \hyperref[regdesc:MINTSTATUS]{mintstatus}
\item \hyperref[regdesc:MCLICBASE_REG]{mclicbase}
\item \hyperref[regdesc:MSCRATCHCSW_REG]{mscratchcsw}
\item \hyperref[regdesc:MSCRATCHCSWL_REG]{mscratchcswl}
\end{itemize}
\end{multicols}



In CLIC mode, the global interrupt enable for machine
mode is still dictated by the \hyperref[fielddesc:MSTATUSMIE]{mstatus.MIE} bit.


Each interrupt can be either vectored or non-vectored depending on the \hyperref[fielddesc:MCLICINTATTR_i_SHV]{clicintattr[i].SHV} bit.

For a non-vectored interrupt, the HP core jumps to the trap handler base address as configured in \hyperref[fielddesc:MTVECBASE]{mtvec.BASE} field.

The base address of the vector table is configured in \hyperref[regdesc:MTVT]{mtvt}.
For a vectored interrupt, the HP core jumps to the 32-bit word address stored at the corresponding index in the vector table.
For example, for the "i"th machine mode interrupt, first the 32-bit address stored at the memory location "mtvt + 4*i" will be loaded, after which the HP core will treat the loaded data as an address and jump to it.

Upon taking an interrupt, the HP core will store the address of the instruction that was interrupted in \hyperref[regdesc:MEPC]{mepc}.
Also, the HP core will update the value of \hyperref[regdesc:MCAUSE]{mcause}
based on the interrupt's ID.


\hypertarget{riscvcpu-clic-intmap}{}
\subsubsection{Interrupt Mapping}
\label{sec:riscvcpu-clic-intmap}

CLIC in the HP core supports 32 external interrupts and 2 core-local interrupts.

The IDs 16-47 are mapped to the external interrupts, while IDs 15 and below are reserved for the core-local interrupts.

\begin{longtable}{ | l | l | }
\caption{CLIC Interrupt Mapping}
\label{tab:riscvcpu-clic-intmap-table}
\\ \hline
\rowcolor{lightgray}
ID & Description \\ \hline
0-2 & NA (Reserved) \\ \hline
3 & CLINT M mode software interrupt \\ \hline
4-6 & NA (Reserved) \\ \hline
7 & CLINT M mode timer interrupt \\ \hline
8-15 & NA (Reserved) \\ \hline
16 & External interrupt 0 \\ \hline
17 & External interrupt 1 \\ \hline
18 & External interrupt 2 \\ \hline
... & ... \\ \hline
45 & External interrupt 29 \\ \hline
46 & External interrupt 30 \\ \hline
47 & External interrupt 31 \\ \hline
\end{longtable}

For detailed description of the core-local interrupt sources, please refer to Section \ref{sec:riscvcpu-clint}.

The mapping of the 32 external interrupts to the interrupt sources in SoC is managed using the Interrupt Matrix. Please refer to Chapter \ref{mod:intmtrx} \textit{\nameref{mod:intmtrx}} for details on how to configure and map to SoC interrupts to CLIC interrupts for the HP core.

\hypertarget{riscvcpu-clic-params}{}
\subsubsection{CLIC Parameters}
\label{sec:riscvcpu-clic-params}

The available features of a CLIC implementation can be discovered using the registers \hyperref[regdesc:MCLICBASE_REG]{mclicbase}
and \hyperref[regdesc:CLICINFO_REG]{clicinfo}:
\begin{itemize}
    \item \hyperref[regdesc:MCLICBASE_REG]{mclicbase} is a read-only CSR which holds the base address of the CLIC memory-mapped registers for machine mode.
    \item \hyperref[regdesc:CLICINFO_REG]{clicinfo} is a memory-mapped read-only register. It consists of the following fields:
    \begin{itemize}
        \item \hyperref[fielddesc:CLICINFO_NUM_INTERRUPTS]{NUM\_INTERRUPTS}: Represents the number of interrupts supported by the current CLIC implementation
        \item \hyperref[fielddesc:CLICINFO_CLICINTCTLBITS]{CLICINTCTLBITS}: Represents the number of programmable higher bits in \hyperref[regdesc:MCLICINTCTL_REG_i]{clicintctl[i]} and \hyperref[fielddesc:MINTTHRESH_TH]{{mintthresh.TH}} for encoding the level and priority of each interrupt
    \end{itemize}
    \item Once the value of the \hyperref[fielddesc:MCLICBASE]{mclicbase} is known, the full address of the  \hyperref[regdesc:CLICINFO_REG]{clicinfo} can be computed by adding the corresponding address offset of 0x0004 (shown in Section \ref{sec:riscvcpu-clic-memmap-reg-summ}) to the \hyperref[fielddesc:MCLICBASE]{mclicbase} value.
\end{itemize}


\hypertarget{riscvcpu-clic-int-level-prio}{}
\subsubsection{Interrupt Level and Priority Encoding}
\label{sec:riscvcpu-clic-int-level-prio}

Table \ref{tab:riscv-clic-level-priority} shows the possible ways that level and priority can be specified for CLIC interrupts, based on the configured values of \hyperref[fielddesc:MCLICCFG_MNLBITS]{mcliccfg.MNLBITS}
and \hyperref[regdesc:MCLICINTCTL_REG_i]{clicintctl[i]} registers, and the fixed value of the \hyperref[fielddesc:CLICINFO_CLICINTCTLBITS]{CLICINTCTLBITS} parameter (0x3 for this CLIC implementation).

\begin{longtable}{ | l | l | l | l | }
\caption{CLIC Level and Priority Encoding with CLICINTCTLBITS=3}
\label{tab:riscv-clic-level-priority}
\\\hline
\rowcolor{lightgray}
mnlbits[3:0] & clicintctl[i][7:0] & Interrupt Levels & Interrupt Priorities \\ \hline
0   & \texttt{ppp.....} & 255 & 31, 63, 95, 127, 159, 191, 223, 255 \\ \hline
1   & \texttt{lpp.....} & 127, 255 & 31, 63, 95, 127 \\ \hline
2   & \texttt{llp.....} & 63, 127, 191, 255 & 31, 63 \\ \hline
3-8 & \texttt{lll.....} & 31, 63, 95, 127, 159, 191, 223, 255 & 31 \\ \hline
\end{longtable}

\hyperref[fielddesc:MCLICCFG_MNLBITS]{mcliccfg.MNLBITS} decides the number of higher bits in the \hyperref[regdesc:MCLICINTCTL_REG_i]{clicintctl[i]} registers which are to be used to encode the level (indicated as \texttt{'l'} in the above table), while the remaining lower bits are used to encode the priority (indicated as \texttt{'p'} in the above table). Since only the higher bits of \hyperref[regdesc:MCLICINTCTL_REG_i]{clicintctl[i]} are writeable as per the \hyperref[fielddesc:CLICINFO_CLICINTCTLBITS]{CLICINTCTLBITS}) parameter, hence the remaining lower bits are tied to 1 (indicated as '.' in the above table).

Level and priority are both used to give preference to one interrupt over the other. A higher level interrupt can preempt a lower level interrupt (assuming that \hyperref[fielddesc:MSTATUSMIE]{mstatus.MIE} is manually enabled inside an interrupt handler for allowing nesting), but in the case of two interrupts with same level and different priorities, the higher priority interrupt cannot preempt the lower priority interrupt while it is being serviced by the CPU. Thus, priority is only used as a tie-breaker to decide between two pending interrupts with the same level. Finally, if both level and priority are same for two interrupts, the one with the higher ID value is given preference.


The CPU internally uses the effective interrupt level threshold, which is defined as the maximum of \hyperref[fielddesc:MINTSTATUS_MIL]{{mintstatus.MIL}} and \hyperref[fielddesc:MINTTHRESH_TH]{{mintthresh.TH}}, to decide if a new pending machine mode interrupt will be allowed to preempt the execution of an ongoing machine mode interrupt handler.

%%%%%%%%%%%%%%%%%%%%%%%%
%%% Register Summary %%%
%%%%%%%%%%%%%%%%%%%%%%%%

\hypertarget{riscvcpu-clic-memmap-reg-summ}{}
\subsubsection{CLIC Memory-Mapped Register Summary}
\label{sec:riscvcpu-clic-memmap-reg-summ}

Listed below are all the CLIC machine mode memory-mapped registers. The address of each register is specified as an offset to the value of  \hyperref[fielddesc:MCLICBASE]{mclicbase}, as well as the absolute address.

Note that the \hyperref[regdesc:MCLICINTIP_REG_i]{clicintip[i]}, \hyperref[regdesc:MCLICINTIE_REG_i]{clicintie[i]}, \hyperref[regdesc:MCLICINTATTR_REG_i]{clicintattr[i]} and \hyperref[regdesc:MCLICINTCTL_REG_i]{clicintctl[i]} registers are all byte sized. Still word and half-word aligned accesses to these registers will be performed atomically. Please note \hyperref[regdesc:MINTTHRESH]{mintthresh} is a machine mode CSR instead of memory mapped register.

\subfile{\modulefiles/01-RISCV-CLIC-M-reg-summ__EN}


%%%%%%%%%%%%%%%%%%%%%%%%%%%%
%%% Register Description %%%
%%%%%%%%%%%%%%%%%%%%%%%%%%%%

\hypertarget{riscvcpu-clic-memmap-reg}{}
\subsubsection{CLIC Memory-Mapped Register Description}
\label{sec:riscvcpu-clic-memmap-reg}

\subfile{\modulefiles/01-RISCV-CLIC-reg__EN}

%%%%%%%%%%%%%
%%% CLINT %%%
%%%%%%%%%%%%%
\newpage
\subsection{Core-Local Interrupts (CLINT)}
\label{sec:riscvcpu-clint}
\subsubsection{Overview}
\label{sec:riscvcpu-clint-overview}
CLINT is the core-local interrupt (CLINT) block that includes software and timer interrupts and corresponding configuration registers. CLINT interrupt sources with IDs as shown below:

\begin{longtable}{ | c | l | }
\caption{Core-Local Interrupt Sources} \label{tab:riscvcpu-clint-table}
\\ \hline
\rowcolor{lightgray}
ID & Description \\ \hline
3 & M mode software interrupt \\ \hline
7 & M mode timer interrupt \\ \hline
\end{longtable}


\subsubsection{Features}
\begin{itemize}
    \item Two local level-type machine mode interrupt sources with programmable priorities
    \item Memory-mapped configuration and status registers
    \item 64-bit timer with interrupt functionality, overflow flags, and three sampling modes
    \item Software interrupt
\end{itemize}

\subsubsection{Software Interrupt}
\label{sec:riscvcpu-clint-software}

An M mode software interrupt is controlled by setting or clearing the memory-mapped register \hyperref[regdesc:CLINTMSIPREG]{MSIP}.

The register bit clicintie[3] must be set for enabling the software interrupt at the core level.

Pending state of this interrupt can be checked by reading the corresponding bit of clicintip register of interrupt 3 in CLIC, that is clicintip[3] bit.


\subsubsection{Timer Counter and Interrupt}
\label{sec:riscvcpu-clint-timer}

The timer counter can be enabled by setting the \hyperref[fielddesc:CLINTMTIMECTLEN]{mtime\_EN} bit in  \hyperref[regdesc:CLINTMTIMECTLREG]{mtimectl}.
\hyperref[regdesc:CLINTMTIMECMPLOREG]{mtimecmplo} and \hyperref[regdesc:CLINTMTIMECMPHIREG]{mtimecmphi} registers are set to the value for comparison with system timer with lower and higher 32 bits respectively.
The CPU provides two local memory-mapped 32-bit wide registers \hyperref[regdesc:CLINTMTIMELOADLOREG]{mtimeloadlo} and \hyperref[regdesc:CLINTMTIMELOADHIREG]{mtimeloadhi}, which has both read/write access, to load the system timer with a specific value.

Interrupt for M mode is asserted when 64-bit value of system timer value exceeds the 64-bit combined value of \hyperref[regdesc:CLINTMTIMECMPLOREG]{mtimecmplo} and \hyperref[regdesc:CLINTMTIMECMPHIREG]{mtimecmphi}.

Pending state of timer interrupt is reflected as IP bit of clicintip register of interrupt 7 in CLIC.

For de-asserting the pending timer interrupt in M mode, either IE bit of clicintie register of interrupt 7 in CLIC has to be cleared or the value of the \hyperref[regdesc:CLINTMTIMECMPLOREG]{mtimecmplo} and \hyperref[regdesc:CLINTMTIMECMPHIREG]{mtimecmphi} register needs to be updated.

%%%%%%%%%%%%%%%%%%%%%%%%
%%% Register Summary %%%
%%%%%%%%%%%%%%%%%%%%%%%%

\hypertarget{riscvcpu-clint-reg-summ}{}
\subsubsection{Register Summary}
\label{sec:riscvcpu-clint-reg-summ}

The addresses in this section are relative to CPU sub-system base address i.e. 0x20000000.


\subfile{\modulefiles/01-RISCV-CLINT-reg-summ__EN}

%%%%%%%%%%%%%%%%%
%%% Registers %%%
%%%%%%%%%%%%%%%%%

\subsubsection{Register Description}

The addresses in this section are relative to CPU sub-system base address i.e. 0x20000000.

\subfile{\modulefiles/01-RISCV-CLINT-reg__EN}

%%%%%%%%%%%%%%%%%%%%%%%%
%%% Memory Protection Unit %%%
%%%%%%%%%%%%%%%%%%%%%%%%
\newpage
\section{Memory Protection Unit}
HP core implements standard RISC-V physical memory protection (PMP) scheme along with custom physical memory attribute checker (PMAC) logic to prevent unauthorized memory access.
%%%%%%%%%%%%%%%%%%%%%%%%
%%% Physical Memory Protection %%%
%%%%%%%%%%%%%%%%%%%%%%%%

\subsection{Standard Physical Memory Protection}
\label{sec:riscvcpu-pmp}
\subsubsection{Overview}
\label{sec:riscvcpu-pmp-overview}
The CPU core includes a physical memory protection (PMP) unit, which is fully compliant with RISC-V Instruction Set Manual, Volume II: Privileged Architecture, Version 1.10. It can be used by software to set memory access privileges (read, write and execute permissions) for required memory regions. In addition to standard PMP checks, the CPU core also implements custom physical memory attributes (PMA) checkers to provide additional permission checks based on pre-defined attributes.

\subsubsection{Features}
\label{sec:riscvcpu-pmp-features}

The PMP unit can be used to restrict access to physical memory. The main features are:

\begin{itemize}
    \item Up to 16 PMP entries
    \item Minimum address granularity of 128 bytes
    \item Three address matching modes: OFF, TOR (Top of Range), and NAPOT (Naturally Aligned Power-of-Two). NA4 (Naturally Aligned 4-Byte Region) address matching mode is not supported
    \item Permission controls for read, write, and execute
    \item Lock function for each PMP entry
    \item Maximum physical space address of 4 GB
\end{itemize}

\subsubsection{Functional Description}
\label{sec:riscvcpu-pmp-func-descr}
Software can program the PMP unit's configuration and address registers in order to contain faults and support secure execution. PMP CSRs can only be programmed in machine mode. Once the PMP unit is enabled by configuring PMP CSRs, write, read and execute permission checks are applied to all the accesses in user mode as per programmed values of the enabled pmpcfg\regindex{n} (\regindex{n}: 0-3) and pmpaddr\regindex{n} (\regindex{n}: 0-15) registers (refer to Section \ref{sec:riscvcpu-pmp-reg-summ}).

By default, PMP grants permission to all accesses in machine mode and revokes permission of all access in user mode. This implies that it is mandatory to program the address range and valid permissions in pmpcfg and pmpaddr registers (refer to Section \ref{sec:riscvcpu-pmp-reg-summ}) for any valid access to pass through in user mode. However, it is not required for machine mode as PMP permits all accesses to go through by default. In cases where PMP checks are also required in machine mode, software can set the lock bit of required PMP entry to enable permission checks on it. Once the lock bit is set, it can only be cleared through CPU reset.

When any instruction is being fetched from a memory region without execute permissions, an exception is generated at processor level and exception cause is set as instruction access fault in \hyperref[regdesc:MCAUSE]{mcause} CSR. Similarly, any load/store access without valid read/write permissions will result in an exception generation with \hyperref[regdesc:MCAUSE]{mcause} updated as load access and store access fault respectively. In case of load/store access faults, violating address is captured in \hyperref[regdesc:MTVAL]{mtval} CSR.

%%%%%%%%%%%%%%%%%%%%%%%%
%%% Register Summary %%%
%%%%%%%%%%%%%%%%%%%%%%%%

\newpage
\hypertarget{riscvcpu-pmp-reg-summ}{}
\subsubsection{Register Summary}
\label{sec:riscvcpu-pmp-reg-summ}

Below is a list of PMP CSRs supported by the CPU. These are only accessible from machine mode.

\subfile{\modulefiles/01-RISCV-PMP-reg-summ__EN}

%%%%%%%%%%%%%%%%%
%%% Registers %%%
%%%%%%%%%%%%%%%%%

\subsubsection{Register Description}
\label{sec:riscvcpu-pmp-reg-desc}

PMP unit implements all pmpcfg0-3 and pmpaddr0-15 CSRs as defined in RISC-V Instruction Set Manual, Volume II: Privileged Architecture, Version 1.10.

\subfile{\modulefiles/01-RISCV-PMP-reg__EN}

%%%%%%%%%%%%%%%%%%%%%%%%%%%%%%%%%%%%%%%%%%%%%%
%%% Physical Memory Attribute (PMA) Checker %%%
%%%%%%%%%%%%%%%%%%%%%%%%%%%%%%%%%%%%%%%%%%%%%%

\newpage
\subsection{Custom Physical Memory Attribute (PMA) Checker}
\label{sec:riscvcpu-pma}
\subsubsection{Overview}
\label{sec:riscvcpu-pma-overview}
CPU core also implements a custom physical memory attributes checker (PMAC) to provide additional permission checks based on pre-defined memory types configured through custom CSRs. Please note that the PMA checks are applied irrespective of the core's privilege mode, i.e., it doesn't differentiate between machine and user mode.


\subsubsection{Features}
\label{sec:riscvcpu-pma-features}

PMAC supports below features:

\begin{itemize}
    \item 16 PMA entries
    \item Configurable memory type for defined memory regions
    \item Minimum address granularity of 128 bytes
    \item Three address matching mode: OFF, TOR (Top of Range) and NAPOT (Naturally Aligned Power-of-Two). NA4 (Naturally Aligned 4-Byte Region) address matching mode is not supported
    \item Permission controls for read, write and execute
    \item Lock function for each PMA entry
    \item Maximum physical space address of 4 GB
    \item Configurable attribute for defined memory regions
\end{itemize}

\subsubsection{Functional Description}
\label{sec:riscvcpu-pma-func-descr}

Software can program the PMAC unit's configuration and address registers in order to avoid faults due to access to invalid memory regions. PMAC CSRs can only be programmed in machine mode. Once enabled, write, read and execute permission checks are applied to all the accesses irrespective of privilege mode as per programmed values of enabled 16 pma\_cfgX and pma\_addrX registers (refer to Section \ref{sec:riscvcpu-pma-reg-summ}). Access to entries marked as invalid memory types will result in fetch fault or load/store fault exception, as the case may be.

Exception generation and handling for PMAC related faults will be handled in similar way to PMP checks. When any instruction is being fetched from a memory region configured as a null or an invalid memory region, an exception is generated at the processor level and the exception cause is set as instruction access fault in \hyperref[regdesc:MCAUSE]{mcause} CSR. Similarly, any load/store access to a null or an invalid memory region will result in an exception generation with \hyperref[regdesc:MCAUSE]{mcause} updated as load access or store access fault respectively. In case of load/store access faults, violating address is captured in \hyperref[regdesc:MTVAL]{mtval} CSR. For the PMAC entries configured as valid memory, the handling is same as for PMP checks.

A lock bit per entry is also provided in case software wants to disable programming of PMAC registers. Once the lock bit in any pma\_cfgX register is set, respective pma\_cfgX and pma\_addrX registers can not be programmed further, unless a CPU reset cycle is applied.

A 4-bit field in PMAC CSRs is also provided to define attributes for memory regions. These bits are not used internally by the CPU core for any purpose. Based on address match, these attributes are provided on load/store interface as side-band signals and are used by cache controller block for its internal operation.

%%%%%%%%%%%%%%%%%%%%%%%%
%%% Register Summary %%%
%%%%%%%%%%%%%%%%%%%%%%%%

\newpage
\hypertarget{riscvcpu-pma-reg-summ}{}
\subsubsection{Register Summary}
\label{sec:riscvcpu-pma-reg-summ}


Below is a list of PMA CSRs supported by the CPU. These are only accessible from machine mode.

\subfile{\modulefiles/01-RISCV-PMA-reg-summ__EN}

%%%%%%%%%%%%%%%%%
%%% Registers %%%
%%%%%%%%%%%%%%%%%

\subsubsection{Register Description}
\label{sec:riscvcpu-pma-reg-desc}

\subfile{\modulefiles/01-RISCV-PMA-reg__EN}

%%%%%%%%%%%%%%%%%%%%%%%%%%%%
%%% Debug & Trace Support %%%
%%%%%%%%%%%%%%%%%%%%%%%%%%%%
\newpage
\section{Debug and Trace Support}

%%%%%%%%%%%%%%%%%%%%%%%%%%%%
%%% Debug %%%
%%%%%%%%%%%%%%%%%%%%%%%%%%%%

\subsection{Debug}
\label{sec:riscvcpu-dbg}

\subsubsection{Overview}
\label{sec:riscvcpu-dbg-overview}

This section describes how to debug software running on HP core. Debug support is provided through standard JTAG pins and complies to the \href{https://riscv.org/wp-content/uploads/2019/03/riscv-debug-release.pdf}{RISC-V External Debug Support Version 0.13.2} specification.

HP core has its own JTAG DTM (debug transport module) in order to connect with its Debug Module.

Figure \ref{fig:riscvcpu-dbg-sys-overview} below shows the main components of External Debug Support.

\begin{figure}[H]
\centering
\includegraphics[width=0.8\textwidth,keepaspectratio]{\modulefiles/figures/debug_bigfonts.png}
\caption{Debug System Overview}\label{fig:riscvcpu-dbg-sys-overview}
\end{figure}


The user interacts with the Debug Host (e.g. laptop), which is running a debugger (e.g. GDB). The debugger communicates with a Debug Translator (e.g. OpenOCD, which may include a hardware driver) to communicate with Debug Transport Hardware (e.g. ESP-Prog adapter). The Debug Transport Hardware connects the Debug Host to the HP core's Debug Transport Module (DTM) through standard JTAG interface. The DTM provides access to the Debug Module (DM) using the Debug Module Interface (DMI).

The DM allows the debugger to halt the HP core. Abstract commands provide access to GPRs (general-purpose registers). The Program Buffer allows the debugger to execute arbitrary code on the core, which allows access to additional CPU core state. Alternatively, additional abstract commands can provide access to additional CPU core state.

HP core contains Trigger Module supporting 3 triggers. When trigger conditions are met, respective core will halt spontaneously and inform the debug module that they have halted.

System bus access block allows memory and peripheral register access without affecting either HP core.


\subsubsection{Features}
\label{sec:riscvcpu-dbg-features}

HP CPU supports the following basic debugging features:

\begin{itemize}
    \item Provides necessary information about the implementation to the debugger
    \item Allows the core to be halted and resumed
    \item Core registers (including CSRs) can be read/written by debugger
    \item Core can be debugged from the first instruction executed after reset
    \item Core can be reset through debugger
    \item Core can be halted on software breakpoint (planted breakpoint instruction)
    \item Hardware single-stepping
    \item Execute arbitrary instructions in the halted core by means of the program buffer. 2-word program buffer is supported
    \item System bus access is supported through word aligned address access
    \item Supports three Hardware Triggers (can be used as breakpoints/watchpoints) as described in Section \ref{sec:riscvcpu-hwt}
\end{itemize}

\subsubsection{Functional Description}
\label{sec:riscvcpu-dbg-func-descr}

As mentioned earlier, Debug Scheme conforms to the \href{https://riscv.org/wp-content/uploads/2019/03/riscv-debug-release.pdf} {RISC-V External Debug Support Version 0.13.2 specification}. Please refer to it for functional operation details.

\subsubsection{JTAG Control}
\label{sec:riscvcpu-dbg-jtag-cntl}

Standard JTAG interface is the only way for DTM to access DM. The hardware provides two JTAG methods: PAD\_to\_JTAG and USB\_to\_JTAG.

\begin{itemize}
    \item {PAD\_to\_JTAG: The JTAG's signal source comes from IO.}
    \item {USB\_to\_JTAG: The JTAG's signal source comes from USB Serial/JTAG Controller.}
\end{itemize}

Which JTAG method to use depends on many factors. The following table shows the configuration method.

\subfile{\modulefiles/tables/JTAG_CNTL__EN}



%%%%%%%%%%%%%%%%%%%%%%%%
%%% Register Summary %%%
%%%%%%%%%%%%%%%%%%%%%%%%

\hypertarget{riscvcpu-dbg-reg-summ}{}
\subsubsection{Register Summary}
\label{sec:riscvcpu-dbg-reg-summ}

Below is the list of Debug CSRs supported by HP core:
\subfile{\modulefiles/01-RISCV-DBG-reg-summ__EN}

All the debug module registers are implemented in conformance to the specification RISC-V External Debug Support Version 0.13.2. Please refer to it for more details.

%%%%%%%%%%%%%%%%%
%%% Registers %%%
%%%%%%%%%%%%%%%%%


\subsubsection{Register Description}
Below are the details of Debug CSRs supported by HP core:

\subfile{\modulefiles/01-RISCV-DBG-reg__EN}

%%%%%%%%%%%%%%%%%%%%%%%%%%%%%%%
%%% Cross-Triggering Support %%%
%%%%%%%%%%%%%%%%%%%%%%%%%%%%%%%





%%%%%%%%%%%%%%%%%%%%%%%%
%%% Register Summary %%%
%%%%%%%%%%%%%%%%%%%%%%%%





%%%%%%%%%%%%%%%%%
%%% Registers %%%
%%%%%%%%%%%%%%%%%


%%%%%%%%%%%%%%%%%%%%%%%%
%%% Hardware Trigger %%%
%%%%%%%%%%%%%%%%%%%%%%%%

\newpage
\subsection{Hardware Trigger}
\label{sec:riscvcpu-hwt}

\subsubsection{Overview}
\label{sec:riscvcpu-hwt-overview}
HP core implements a hardware trigger block to provide breakpoint and watchpoint capability for debugging. It conforms to the \href{https://riscv.org/wp-content/uploads/2019/03/riscv-debug-release.pdf} {RISC-V External Debug Support Version 0.13.2 specification}. Please refer to it for functional operation details.
\subsubsection{Features}
\label{sec:riscvcpu-hwt-features}

HP core trigger block support following features:

\begin{itemize}
    \item Three independent trigger units, each of which can be configured for matching the address of the program counter or load-store accesses
    \item Preempting execution by causing a breakpoint exception
    \item Halting execution and transferring control to debuggers
    \item Support for NAPOT (naturally aligned power-of-two regions) address encoding
\end{itemize}

\subsubsection{Functional Description}
\label{sec:riscvcpu-hwt-func-descr}

The Hardware Trigger module provides seven CSRs, which are listed under Section \hyperref[sec:riscvcpu-hwt-reg-summ]{register summary}. Among these, \hyperref[regdesc:TDATA1]{tdata1} and \hyperref[regdesc:TDATA2]{tdata2} are abstract CSRs, which means they are shadow registers for accessing internal registers for each of the three trigger units, one at a time.

To choose a particular trigger unit write the index (0-2) of that unit into \hyperref[regdesc:TSELECT]{tselect} CSR. When tselect is written with a valid index, the abstract CSRs \hyperref[regdesc:TDATA1]{tdata1} and \hyperref[regdesc:TDATA2]{tdata2} are automatically mapped to reflect internal registers of that trigger unit. Each trigger unit has two internal registers, namely \hyperref[regdesc:MCONTROL]{mcontrol} and \hyperref[regdesc:MADDRESS]{maddress}, which are mapped to tdata1 and tdata2, respectively.

Writing larger than allowed indexes to tselect will clip the written value to the largest valid index, which can be read back. This property may be used for enumerating the number of available triggers during initialization or when using a debugger.

Since software or debugger may need to know the type of the selected trigger to correctly interpret tdata1 and tdata2, the 4 bits (31-28) of tdata1 encodes the type of the selected trigger. This type field is read-only and always provides a value of 0x2 for every trigger, which stands for match type trigger, hence, it is inferred that tdata1 and tdata2 are to be interpreted as mcontrol and maddress. The information regarding other possible values can be found in the specification, RISC-V External Debug Support Version 0.13.2, but this trigger module only supports type 0x2.

Once a trigger unit has been chosen by writing its index to tselect, it will become possible to configure it by setting the appropriate bits in mcontrol CSR (tdata1) and writing the target address to maddress CSR (tdata2).

Each trigger unit can be configured to either cause breakpoint exception or enter debug mode, by writing to the action field of mcontrol. This bit can only be written from debugger, thus by default a trigger, if enabled, will cause breakpoint exception.

mcontrol for each trigger unit has a \hyperref[fielddesc:MCONTROLHIT]{hit} bit which may be read, after CPU halts or enters exception, to find out if this was the trigger unit that fired. This bit is set as soon as the corresponding trigger fires, but it has to be manually cleared before resuming operation. Although, failing to clear it does not affect normal execution in any way.

Each trigger unit only supports match on address, although this address could either be that of a load/store access or the virtual address of an instruction. The address and size of a region are specified by writing to maddress (tdata2) CSR for the selected trigger unit. Larger than 1 byte region sizes are specified through NAPOT (naturally aligned power-of-two) encoding (see Table \ref{tab:riscvcpu-hwt-table}) and enabled by setting match bit in mcontrol. Note that for NAPOT encoded addresses, by definition, the start address is constrained to be aligned to (i.e. an integer multiple of) the region size. The maximum size supported is 128 bytes.

\begin{longtable}{ | c | c | r | }
\caption{NAPOT encoding for maddress} \label{tab:riscvcpu-hwt-table}
\\ \hline
\rowcolor{lightgray}
maddress(31-0) & Start Address & Size (bytes) \\ \hline
aaa...aaaaaaaaa0 & aaa...aaaaaaaaa0 & $2$ \\ \hline
aaa...aaaaaaaa01 & aaa...aaaaaaaa00 & $4$ \\ \hline
aaa...aaaaaaa011 & aaa...aaaaaaa000 & $8$ \\ \hline
aaa...aaaaaa0111 & aaa...aaaaaa0000 & $16$ \\ \hline
\multicolumn{3}{c}{....} \\ \hline
aaa...0111111111 & aaa...aaa0000000 & $2^{8}$ \\ \hline
\end{longtable}

\hyperref[regdesc:TCONTROL]{tcontrol} CSR is common to all trigger units. It is used for preventing triggers from causing repeated exceptions in machine mode while execution is happening inside a trap handler. This also disables breakpoint exceptions inside ISRs by default, although, it is possible to manually enable this right before entering an ISR, for debugging purposes. This CSR is not relevant if a trigger is configured to enter debug mode.

\subsubsection{Trigger Execution Flow}
\label{sec:riscvcpu-hwt-suggest-op}

When hart is halted and enters debug mode due to the firing of a trigger (\hyperref[fielddesc:MCONTROLACTION]{action} = 1):
\begin{itemize}
    \item \hyperref[regdesc:DPC]{dpc} is set to current PC (in decode stage)
    \item cause field in \hyperref[regdesc:DCSR]{dcsr} is set to 2, which means halt due to trigger
    \item \hyperref[fielddesc:MCONTROLHIT]{hit} bit is set to 1, corresponding to the trigger(s) which fired
\end{itemize}

When hart goes into trap due to the firing of a trigger (\hyperref[fielddesc:MCONTROLACTION]{action} = 0) :
\begin{itemize}
    \item \hyperref[regdesc:MEPC]{mepc} is set to current PC (in decode stage)
    \item \hyperref[regdesc:MCAUSE]{mcause} is set to 3, which means breakpoint exception
    \item \hyperref[fielddesc:TCONTROLMPTE]{mpte} is set to the value in \hyperref[fielddesc:TCONTROLMTE]{mte} right before trap
    \item \hyperref[fielddesc:TCONTROLMTE]{mte} is set to 0
    \item \hyperref[fielddesc:MCONTROLHIT]{hit} bit is set to 1, corresponding to the trigger(s) which fired
\end{itemize}

\textit{Note: If two different triggers fire at the same time, one with action = 0 and another with action = 1, then hart is halted and enters debug mode.}

%%%%%%%%%%%%%%%%%%%%%%%%
%%% Register Summary %%%
%%%%%%%%%%%%%%%%%%%%%%%%

\hypertarget{riscvcpu-hwt-reg-summ}{}
\subsubsection{Register Summary}
\label{sec:riscvcpu-hwt-reg-summ}


Below is a list of Trigger Module CSRs supported by the CPU. These are only accessible from machine mode.

\subfile{\modulefiles/01-RISCV-HWT-reg-summ__EN}

%%%%%%%%%%%%%%%%%
%%% Registers %%%
%%%%%%%%%%%%%%%%%

\subsubsection{Register Description}

\subfile{\modulefiles/01-RISCV-HWT-reg__EN}

%%%%%%%%%%%%
%%% Trace %%%
%%%%%%%%%%%%

\newpage
\subsection{Trace}
\label{sec:riscvcpu-trace}

\subsubsection{Overview}
\label{sec:riscvcpu-trace-overview}
In order to support non-intrusive software debug, the CPU core provides an instruction trace interface which provides relevant information for offline debug purpose. This interface provides relevant information to Trace Encoder block, which compresses the information and stores in memory allocated for it. Software decoders can read this information from trace memory without interrupting the CPU core and re-generate the actual program execution by the CPU core.


\subsubsection{Features}
\label{sec:riscvcpu-trace-features}

The CPU core supports instruction trace feature and provides below information to Trace Encoder as mandated in RISC-V Processor Trace Version 1.0:

\begin{itemize}
    \item Number of instructions being retired.
    \item Occurrence of exception and interrupt along with cause and trap values.
    \item Current privilege level of hart.
    \item Instruction type of retired instructions for jumps, branches and return.
    \item Instruction address for instructions retired before and after program counter changes.
    \item Trace enabling on trigger hit as per \hyperref[regdesc:MCONTROL]{action} bit in \hyperref[regdesc:MCONTROL]{mcontrol}.
\end{itemize}


\subsubsection{Functional Description}
\label{sec:riscvcpu-trace-func-descr}

HP core implements mandatory instruction delta tracing, also known as branch tracing. It works by tracking execution from a known start address by sending information about deltas taken by a program. Deltas are typically introduced by jump, call, return, branch type instructions and also by interrupts and exceptions. All such deltas along with additional details about cause and actual instructions/addresses are communicated over high bandwidth instruction trace interface output from the core. Trace Encoder operates on the information on this trace interface and compresses the information for storage in memory for offline debug by a decoder. More information about the encoding is available in Chapter \ref{mod:riscvtracenc} \textit{\nameref{mod:riscvtracenc}}.

The core does not have any internal registers to provide control over instruction trace interface. All register controls are available in \ref{mod:riscvtracenc} \textit{\nameref{mod:riscvtracenc}} block.

%%%%%%%%%%%%
%%% Performance Enhancements %%%
%%%%%%%%%%%%
\newpage
\section{Performance}
\label{sec:riscvcpu-perf-enhancements}

%%%%%%%%%%%%
%%% Branch Prediction %%%
%%%%%%%%%%%%

\subsection{Branch Prediction}
\label{sec:riscvcpu-perf-bp}

\subsubsection{Overview}
\label{sec:riscvcpu-perf-bp-overview}
A core does branch prediction in terms of branch is taken or not taken and the target address of conditional/unconditional branch type instruction prior to its execution. Based on the prediction, the core starts fetching from the corresponding PC. Then, it validates the prediction on the execution of the instruction. If the predication turns out to be incorrect, the core flushes the pipeline and start fetching from correct PC. Branch prediction improves the throughput by allowing the core to keep fetching instructions without waiting for the completion of earlier branch instruction. The details about the prediction technique used by the core are described in the following section.

\subsubsection{Features}
\label{sec:riscvcpu-perf-bp-features}
\begin{itemize}
    \item Prediction of both results and target address of conditional/unconditional branch type instructions
    \item Branch target prediction of instructions including BEQ, BNE, BLT, BLTU, BGE, BGEU, C.BEQZ, C.BNEZ, JAL, J, CJ
    \item Conditional branch prediction includes instructions such as: BEQ, BNE, BLT, BLTU, BGE, BGEU, C.BEQZ, C.BNEZ
    \item Uses 2-bit saturating counter for more accuracy
    \item Uses global history of conditional branches for prediction which improves performance
    \item Target prediction in one clock cycle for unconditional branch, that is, 0 cycle penalty when prediction or speculation is correct
\end{itemize}

\subsubsection{Functional Description}
\label{sec:riscvcpu-perf-bp-func-descr}

The core uses Branch History Table (BHT) along with Branch Target Buffer (BTB) for prediction of branch and jump type instructions.
A 512x16-bit SRAM memory used for maintaining Branch History Table for HP core. During prediction, it is indexed by Virtual Branch History Register (VBHR). During update, after branch resolves, the BHT is indexed using Global Branch History Register (GBHR), which records the outcomes of the last 13 executed branch instructions. VBHR contains the results of the last 12 branch-type instructions and predicted results of the current branch instruction. Thus, both registers hold the trajectory of previously executed branches. Based on 13 bits of VBHR and GBHR, 2 bits of saturation counter value stored in BHT are either read or written respectively.

As shown in Figure \ref{fig:riscvcpu-bp_sat-scheme}, 2-bit saturation counter indicates if a branch is:

\begin{itemize}
    \item weakly not taken (00)
    \item strongly not taken (01)
    \item weakly taken (10)
    \item strongly taken (11)
\end{itemize}

If the prediction result counter is in any not-taken state, the branch is predicted as not taken. Similarly, if the prediction result counter is in any taken state, the branch is predicted as taken. When the actual branch outcome is known, the BHT entry is updated based on the following 2-bit saturating counter transition logic.

BTB stores the 32-bit branch target address. It can hold a maximum of 16 entries of target addresses. It is indexed by the last 16 bits of the PC of conditional/unconditional branch-type instruction. Thus, the core predicts the result of branch-type instructions based on BHT and jumps to the target obtained from BTB if the branch is predicted as taken. For unconditional branches, the core directly jumps to the target obtained from BTB, as it does not need prediction in this case. When the actual result and the target of branch-type instructions are obtained during execution, the core takes a corrective jump by flushing the pipeline in case of misprediction. BTB and BHT entries are corrected accordingly.

Figure \ref{fig:riscvcpu-bp_sat-scheme} shows 2-bit saturation scheme details used in HP core.
\begin{figure}[h]
\centering
\includegraphics[width=0.8\textwidth,keepaspectratio]{\modulefiles/figures/sat_counter-3.png}
\caption{Two-bit Saturation Counter}\label{fig:riscvcpu-bp_sat-scheme}
\end{figure}

%%%%%%%%%%%%
%%% RAS  %%%
%%%%%%%%%%%%

\subsection{RAS}
\label{sec:riscvcpu-perf-ras}

\subsubsection{Overview}
\label{sec:riscvcpu-perf-ras-overview}
Return Address Stack (RAS) is a stack maintained by the core to push the return addresses of jump and link (JAL/JALR) type instructions. Thus, the core can jump to the return address as soon as the RET instruction is pre-decoded, enhancing the performance of function call return.

\subsubsection{Features}
\label{sec:riscvcpu-perf-ras-features}

\begin{itemize}
    \item Increased core performance by avoiding pipeline stall on return of function call
    \item Able to store up to four levels of function call nesting
\end{itemize}


\subsubsection{Functional Description}
\label{sec:riscvcpu-perf-ras-func-descr}
When JAL/JALR/C.JALR instruction is pre-decoded by the core, the 32-bit PC of the next instruction is pushed to RAS.

The address is popped out from RAS when the return type instruction (JALR, C.JR, C.JALR) is pre-decoded by the core. The core then starts fetching from the popped address. On execution of the return instruction, if the popped address does not match the address stored in the link register, then the core flushes the pipeline and starts to fetch from the address in the link register. RAS can store a maximum of four entries of return addresses, that is, four levels of function call nesting are supported.

RAS prediction depends on the following conditions:

\begin{longtable}{ | c | c | c | }
\hline\rowcolor{lightgray}
Target Register (Rd) & Source Register (Rs) & RAS action \\ \hline
\endfirsthead
\hline
rd != x1 & rs1 != x1 & No action \\ \hline
rd != x1 & rs1 == x1 & POP \\ \hline
rd == x1 & rs1 != x1 & PUSH \\ \hline
rd == x1 & rs1 == x1 & PUSH and POP \\ \hline
\end{longtable}

 \subsection{Control Status Register for Performance Configuration}
\label{sec:riscvcpu-perf-csr}
\begin{register}{H}{mhcr}{0x7C1}\label{regdesc:MHCR}
\regfield{(reserved)}{19}{13}{{0x00000}}%
\regfield{BTB}{1}{12}{{0}}%
\regfield{(reserved)}{6}{6}{{0x00}}%
\regfield{BPE}{1}{5}{{0}}%
\regfield{RS}{1}{4}{{0}}%
\regfield{(reserved)}{4}{0}{{0x0}}%
\reglabel{Reset}\regnewline%
\begin{regdesc}\begin{reglist}
\label{fielddesc:MHCRRS}\item [RS] Configures whether to enable the return address stack logic for prediction. \\
0: Disable \\
1: Enable \\
(R/W)
\label{fielddesc:MHCRBPE}\item [BPE] Configures whether to enable conditional branch prediction.\\
0: Disable \\
1: Enable \\
(R/W)
\label{fielddesc:MHCRBTB}\item [BTB] Configures whether to enable branch target buffer.\\
0: Disable \\
1: Enable \\
(R/W)
\end{reglist}\end{regdesc}
\end{register}
%%%%%%%%%%%%%%%%%%%%%%%%%%%%%%%
%%% Custom Features %%%
%%%%%%%%%%%%%%%%%%%%%%%%%%%%%%%
\newpage
\section{Custom Features}

%%%%%%%%%%%%%%%%%%%%%%%%%
%%% Bus Error Response %%%
%%%%%%%%%%%%%%%%%%%%%%%%%

\subsection{Bus Error Response}
\label{sec:riscvcpu-bus-err-resp}
\subsubsection{Overview}
\label{sec:riscvcpu-buserr-overview}

In the HP core, load/store bus errors are handled as asynchronous exceptions. As a result, the CPU may have executed some instructions following the load/store instruction that caused the exception before it is detected.

Due to the asynchronous nature of these exceptions, handling them requires more context than that available via the standard CSRs \hyperref[regdesc:MCAUSE]{mcause}, \hyperref[regdesc:MEPC]{mepc} and \hyperref[regdesc:MTVAL]{mtval}, which are otherwise sufficient for synchronous exceptions. Therefore, additional CSRs have been provided to allow software to gain all the necessary context required to handle and resume a program that has encountered an asynchronous bus-error exception.

In the following sections, these bus-error-specific CSRs are described and the recommended approach is to gain context about a bus-error exception.

\subsubsection{Functional Description}
\label{sec:riscvcpu-buserr-func-descr}

There are two ways to handle load/store bus errors:

\begin{itemize}
    \item Synchronous bus-errors mode: \\
    When this mode is enabled, (by setting \hyperref[fielddesc:MHINT_SBE]{mhint.SBE} bit) bus-error exceptions can be handled just like other synchronous exceptions. Therefore, execution traps at the exact PC corresponding to the bus error causing load/store instruction, and \hyperref[regdesc:MEPC]{mepc} holds the same PC.\\
    The side effect of enabling this mode is that it increases the latency of all load/store instructions.
    \item Asynchronous bus-errors mode (default):
    \begin{itemize}
        \item For handling bus errors in asynchronous mode, it is recommended (although not necessary) to clear the \hyperref[fielddesc:MEXSTATUS_NMFT]{mexstatus.NMFT} bit at startup.\\
        \item In asynchronous mode, bus errors cause the CPU to enter trap handler with MCAUSE value 5 or 7 (depending upon the type of operation which caused the error), and \hyperref[regdesc:MEPC]{mepc} will hold the PC of the last instruction that was about to be executed when CPU detected the bus error(s) due to (one or more) previously executed load/store instruction(s).\\
         \item Upon entering the exception handler for cause 5 or 7, the \hyperref[fielddesc:MEXSTATUS_BUS_ERR]{mexstatus.BUS\_ERR} bit must be checked to confirm if it is due to load/store bus-error. If so, this bit must be cleared to acknowledge.\\
         \item Next, bits \hyperref[fielddesc:LDPC0_VLD]{ldpc0.vld}, \hyperref[fielddesc:LDPC1_VLD]{ldpc1.vld}, \hyperref[fielddesc:STPC0_VLD]{stpc0.vld}, \hyperref[fielddesc:STPC1_VLD]{stpc1.vld} and \hyperref[fielddesc:STPC2_VLD]{stpc2.vld} must be checked to confirm if they are valid. If so, these must be cleared, and the corresponding address fields (\hyperref[fielddesc:LDPC0_ADDR]{ldpc0.addr}, \hyperref[fielddesc:LDPC1_ADDR]{ldpc1.addr}, \hyperref[fielddesc:STPC0_ADDR]{stpc0.addr}, \hyperref[fielddesc:STPC1_ADDR]{stpc1.addr}, and \hyperref[fielddesc:STPC2_ADDR]{stpc2.addr}) can be extracted from these CSRs to get the program counters at which these faulting load/store instructions were executed. The CSRs \hyperref[regdesc:LDTVAL0]{ldtval0}, \hyperref[regdesc:LDTVAL1]{ldtval1}, \hyperref[regdesc:STTVAL0]{sttval0}, \hyperref[regdesc:STTVAL1]{sttval1} and \hyperref[regdesc:STTVAL2]{sttval2} hold the memory access addresses corresponding to the respective loads/stores. Note that the lower indexed CSRs hold the newest bus errors, while the higher indexed CSRs hold the older bus errors.\\
         \item Note that, even if the cause is either 5 or 7, both load and store bus error CSRs must be checked as sometimes there are more than one load/store instructions (or misaligned accesses which are split into multiple loads/stores) that have been executed and are in flight. Prior to entering the trap handler for the first encountered bus error, the CPU will allow all pending memory operations to finish and record any bus errors caused by them, thus it is important to check all the CSRs to make sure that all such exceptions were caught and handled.\\
         \item Also note that recovering a normal program execution after encountering a bus-error exception is not as simple as for synchronous cases, since program execution may have proceeded quite far ahead of the bus-error-causing instructions.
    \end{itemize}
\end{itemize}

\subsubsection{Control Status Registers}
\label{sec:riscvcpu-buserr-csr}

\begin{longtable}{ | p{4.0cm} | p{8.4cm} | c | c| }
\hline\rowcolor{lightgray}
Name & Description & Address & Access \\ \hline
\endfirsthead
\hline
\rowcolor{lightgray}
Name & Description & Address & Access  \\
\hline
\endhead
\endfoot

\hyperref[regdesc:MEXSTATUS]{mexstatus} & Machine extension control and status register. & 0x7E1 & R/W \\\hline
\hyperref[regdesc:MHINT]{mhint} & Machine implicit operation register & 0x7C5 & R/W \\\hline
\hyperref[regdesc:LDPC0]{ldpc0} & Load bus error PC register 0 & 0xBE0 & R/W \\\hline
\hyperref[regdesc:LDPC1]{ldpc1} & Load bus error PC register 1 & 0xBE1 & R/W \\\hline
\hyperref[regdesc:LDTVAL0]{ldtval0} & Load bus error access address register 0 & 0xBE8 & R/W \\\hline
\hyperref[regdesc:LDTVAL1]{ldtval1} & Load bus error access address register 1 & 0xBE9 & R/W \\\hline
\hyperref[regdesc:STPC0]{stpc0} & Store bus error PC register 0 & 0xBF0 & R/W \\\hline
\hyperref[regdesc:STPC1]{stpc1} & Store bus error PC register 1 & 0xBF1 & R/W \\\hline
\hyperref[regdesc:STPC2]{stpc2} & Store bus error PC register 2 & 0xBF2 & R/W \\\hline
\hyperref[regdesc:STTVAL0]{sttval0} & Store bus error access address register 0 & 0xBF8 & R/W \\\hline
\hyperref[regdesc:STTVAL1]{sttval1} & Store bus error access address register 1 & 0xBF9 & R/W \\\hline
\hyperref[regdesc:STTVAL2]{sttval2} & Store bus error access address register 2 & 0xBFA & R/W \\\hline
\end{longtable}

%%%%%%%%%%%%%%%%%%%%%%%%%%%%%%%
%%% Dedicated IO %%%
%%%%%%%%%%%%%%%%%%%%%%%%%%%%%%%
\newpage
\subsection{Dedicated IO}
\label{sec:riscvcpu-fastio}
\subsubsection{Overview}
\label{sec:riscvcpu-fastio-overview}

Normally, GPIOs are an APB peripheral, which means that changes to outputs and reads from inputs can get stuck in write buffers or behind other transfers. They are slower because the APB bus runs at a lower speed than the CPU core.
As an alternative, the CPU core implements I/O processor-specific CPU registers (CSRs), which are directly connected to the GPIO matrix or IO pads. These registers are single-instruction accessible, allowing rapid control of the IO pads.

\subsubsection{Features}
\label{sec:riscvcpu-fastio-features}
 \begin{itemize}
    \item 8 dedicated IOs directly mapped on GPIOs
    \item No latency for driving output ports
    \item Two CPU cycle latency for sensing input values
\end{itemize}

\subsubsection{Functional Description}
\label{sec:riscvcpu-fastio-func-descr}

The CPU core has a set of 8 inputs and outputs (pin value + pin output enable). These input and output ports are directly connected to the GPIO matrix, through which they can be mapped on top-level pads. Please refer to Chapter \ref{mod:iomuxgpio} \textit{\nameref{mod:iomuxgpio}} for more details.

The CPU implements three custom CSRs:
\begin{itemize}
    \item GPIO\_IN is read-only and reflects the input value.
    \item GPIO\_OUT is R/W and reflects the output value for the GPIOs.
    \item GPIO\_OEN is R/W and reflects the output enable state for the GPIOs. It controls the pad direction. Programming high would mean the pad should be configured in output mode. Programming low means it should be configured in input mode.
\end{itemize}

%%%%%%%%%%%%%%%%%%%%%%%%
%%% Register Summary %%%
%%%%%%%%%%%%%%%%%%%%%%%%

\newpage
\hypertarget{riscvcpu-fastio-reg-summ}{}
\subsubsection{Register Summary}
\label{sec:riscvcpu-fastio-reg-summ}


Below is a list of custom dedicated IO CSRs implemented inside the core.

\subfile{\modulefiles/01-RISCV-FASTIO-reg-summ__EN}

%%%%%%%%%%%%%%%%%
%%% Registers %%%
%%%%%%%%%%%%%%%%%

\subsubsection{Register Description}
\label{sec:riscvcpu-fastio-reg-desc}
\subfile{\modulefiles/01-RISCV-FASTIO-reg__EN}

%%%%%%%%%%%%%%%%%%%%%%%%
%%% RunStall %%%
%%%%%%%%%%%%%%%%%%%%%%%%

\newpage
\subsection{RunStall Support}
\label{sec:riscvcpu-custom-runstall}
\subsubsection{Overview}
\label{sec:riscvcpu-custom-runstall-overview}

The RunStall signal allows an external agent to stall the processor and shut off most of the clock tree to save operating power.
The RunStall input allows external logic to stall the processor's internal pipeline. Any in-progress instructions in the pipeline will complete after the assertion of the RunStall input line. After the pipeline is stalled, the core outputs a handshake signal "corestalled" to indicate successful stalling of the CPU pipeline.


\subsubsection{Functional Description}
\label{sec:riscvcpu-custom-runstall-feature-func-descr}

The RunStall input can be used in the following two scenarios:
\begin{itemize}
    \item As a mechanism to initialize instruction and data RAMs after a reset. If the RunStall input is asserted during reset and remains asserted after the reset is removed, the processor will be in a stalled state. Instruction and data RAMs can be initialized and after these RAMs have been initialized, the RunStall can be released, allowing the processor to start code execution. This mechanism allows the processor’s reset vector to point to an address in local instruction RAM, and the reset code can be written to the instruction RAM by an external agent before the processor starts executing code.
    \item To freeze the processor’s internal pipeline. Assertion of RunStall reduces the processor’s active power dissipation without using or requiring interrupts and the WAITI option. However, the power savings resulting from use of the RunStall signal will not be as great as for those using WAITI, because the RunStall logic allows some portions of the processor to be active even though the pipeline is stalled.
\end{itemize}

Figure \ref{fig:riscvcpu-runstall} shows the timing diagram for RunStall and CoreStalled signals and CPU core state transition.
\begin{figure}[h]
\centering
\includegraphics[width=0.8\textwidth,keepaspectratio]{\modulefiles/figures/RunStall.png}
\caption{CPU Core State Transition Using RunStall }\label{fig:riscvcpu-runstall}
\end{figure}

\subsubsection{Register Summary}
\label{sec:riscvcpu-custom-runstall-reg-summ}

The \hyperref[fielddesc:TRACESTALLENA]{TRACE\_STALL\_ENA} register field in Chapter \ref{mod:riscvtracenc} \textit{\nameref{mod:riscvtracenc}} can be used to assert RunStall to HP core.

%%%%%%%%%%%%%%%%%%%%%%%%
%%% Debug Assist Information %%%
%%%%%%%%%%%%%%%%%%%%%%%%

\subsection{Debug Assist Information}
\label{sec:riscvcpu-dbgassist-features}
The HP CPU core complex provides additional debug information to the Debug Assistant module in \chipname{}. For details, please refer to Chapter \ref{mod:dbgassist} \textit{\nameref{mod:dbgassist}}.





%%%%%%%%%%%%%%%%%%%%%%%%
%%% CPU Lockup %%%
%%%%%%%%%%%%%%%%%%%%%%%%

\subsection{Core Lock-up}
\label{sec:riscvcpu-custom-lockup-features}
\subsubsection{Overview}
\label{sec:riscvcpu-custom-lockup-overview}
Currently, there is no exception enable register defined in the RISC-V programming model, and when returning from the exception service program to the normal program flow, the exception vector number in the MCAUSE register is not cleared, so software developers cannot easily know whether the CPU is currently processing an exception or running a normal program track by the registers defined in the programming model. To address this concern, a custom LOCKUP field has been implemented. It is readable through the custom CSR.

\subsubsection{Functional Description}
\label{sec:riscvcpu-custom-lockup-feature-func-descr}
When the CPU responds to an exception, it will determine whether the current program flow is in the exception service program (through EXPT\_VLD in MEXSTATUS field). If the CPU is processing an exception and triggers a new exception before returning from the exception service program, the CPU will be locked. When CPU is locked:

\begin{itemize}
    \item the lock-up signal is set high to inform the SoC that the CPU is in a locked state
    \item the CPU stops instruction fetch and execution
    \item the CPU keeps the PC at the instruction PC that triggers the CPU lock
    \item the LOCKUP field in the MEXSTATUS register is set to 1
\end{itemize}

In a locked-up state, a debugger can still break in and execute debug commands. It is recommended to reset the core when a lock state is detected.

There are some exceptions to the lock-up rule. Exceptions caused by ECALL or EBREAK, when nested in any exception type, will not trigger CPU lock-up.

\subsubsection{Register Summary}
\label{sec:riscvcpu-custom-lockup-reg-summ}

The \hyperref[fielddesc:MEXSTATUS_LOCKUP]{mexstatus.LOCKUP} bit provides lock-up status of individual core.

\subsubsection{Register Description}
\label{sec:riscvcpu-custom-lockup-reg-desc}
Please refer to \hyperref[regdesc:MEXSTATUS]{mexstatus} CSR.

\end{document}
